% problem-000409 2 配位结构与化学计算
\section*{2. 配位结构与化学计算}
\textit{下列为分问。}

\subsection*{2-1}
$\mathrm { B i _ { 2 } C l _ { 8 } } ^ { 2 \mathrm { - } }$ −离⼦中铋原⼦的配位数为5，配体呈四⻆锥型分布，画出该离⼦的结构并指出Bi原⼦的杂化轨道类型。

\subsection*{2-2}
在液氨中， $E ^ { \ominus } \mathrm { ( N a ^ { + } / N a ) } = - 1 . 8 9 \mathrm { V } , E ^ { \ominus } \mathrm { ( M g ^ { 2 + } / M g ) } = - 1 . 7 4 \mathrm { V }$ ，但可以发⽣ $\mathrm { M g }$ 置换Na的反应： $\mathrm { M g + }$ $2 \mathrm { N a I }  \mathrm { M g I } _ { 2 } + 2 \mathrm { N a }$ ，指出原因。

\subsection*{2-3}
将Pb加到氨基钠的液氨溶液中，先⽣成⽩⾊沉淀 $\mathrm { N a } _ { 4 } \mathrm { P b }$ ，随后转化为 $\mathrm { N a } _ { 4 } \mathrm { P b } _ { 9 }$ （绿⾊）⽽溶解。在此溶液中插⼊两块铅电极，通直流电，当1.0 mol电⼦通过电解槽时，在哪个电极（阴极或阳极）上沉积出铅？写出沉积铅的量。

\subsection*{2-4}
下图是某⾦属氧化物的晶体结构⽰意图。图中，⼩球代表⾦属原⼦，⼤球代表氧原⼦，细线框出其晶胞。

（图：）

\subsection*{2-4}
-1 写出⾦属原⼦的配位数(�)和氧原⼦的配位数(�)。

\subsection*{2-4}
-2 写出晶胞中⾦属原⼦数(�)和氧原⼦数(�)。

\subsection*{2-4}
-3 写出该⾦属氧化物的化学式（⾦属⽤M表⽰）。

\subsection*{2-5}
向含 $\mathrm { { [ c i s { - } C o ( N H _ { 3 } ) _ { 4 } ( H _ { 2 } O ) _ { 2 } ] } \mathrm { { ; } } }$ 3+的溶液中加⼊氨⽔，析出含 $\{ \mathrm { C o } [ \mathrm { C o } ( \mathrm { N H _ { 3 } } ) _ { 4 } ( \mathrm { O H } ) _ { 2 } ]$ _{3}}6+的难溶盐。

$\{ \mathrm { C o } [ \mathrm { C o } ( \mathrm { N H _ { 3 } } ) _ { 4 } ( \mathrm { O H } ) _ { 2 } ] _ { 3 } \} ^ { 6 + }$ 是以羟基为桥键的多核络离⼦，具有⼿性。画出其结构。

\subsection*{2-6}
向 $\mathrm { K _ { 2 } C r _ { 2 } O _ {  } }$ _{7} 和 $\mathrm { N a C l }$ 的混合物中加⼊浓硫酸制得化合物 X（154.9 g mol−1）。X 为暗红⾊液体，沸点$1 1 7 \ ^ { \circ } \mathrm { C } ,$ ，有强刺激性臭味，遇⽔冒⽩烟，遇硫燃烧。X分⼦有两个相互垂直的镜⾯，两镜⾯的交线为⼆重旋转轴。写出X的化学式并画出其结构式。

\subsection*{2-7}
实验得到⼀种含钯化合物 $\mathrm { P d } [ \mathrm { C } _ { x } \mathrm { H } _ { y } \mathrm { N } _ { z } ] ( \mathrm { C l O } _ { 4 } ) _ { 2 }$ ，该化合物中C和H的质量分数分别为30.15%和

$5 . 0 6 \%$ 。将此化合物转化为硫氰酸盐 $\mathrm { P d } [ \mathrm { C } _ { x } \mathrm { H } _ { y } \mathrm { N } _ { z } ] ( \mathrm { S C N } )$ _{2}， 则C 和H 的质量分数分别为40.46%和5.94%。通过计算确定 $\mathrm { P d } [ \mathrm { C } _ { x } \mathrm { H } _ { y } \mathrm { N } _ { z } ] ( \mathrm { C l O } _ { 4 } )$ _{2}的组成。

\subsection*{2-8}
甲烷在汽⻋发动机中平稳、完全燃烧是保证汽⻋安全和⾼能效的关键。甲烷与空⽓按⼀定⽐例混合，氧⽓的利⽤率为85%，计算汽⻋尾⽓中O_{2}、 $\mathrm { C O } _ { 2 } .$ $_ \mathrm { H _ { 2 } O }$ 和 $\mathrm { N } _ { 2 }$ 的体积⽐（空⽓中 $\mathrm { O _ { 2 } }$ 和 $\mathrm { N } _ { 2 }$ 体积⽐按21:79计；设尾⽓中 $\mathrm { { C O } _ { 2 } }$ 的体积为1）。
